% Section 9.4, from lectures/L09/appendix/c_exercises.md. The note under the lecture's intro, which
% exercise is the Kontroll, is folded into the aside, as in Chapter 1. Exercise 9 uses the whole of
% long_delay.asm, lines outside the delay included, so the program is printed there in full.

\section{Övningar}
\label{c9:sec:exercises}\label{c9:app:c}

\begin{aside}
\textbf{Så kontrollerar du ditt arbete.} Rita flödesplanen innan du skriver programmet, och
förutsäg resultatet innan du simulerar. Övningarna märkta \textbf{Program} och \textbf{Kontroll}
görs i Microchip Studio; tidsfördröjningar mäts med stoppuret enligt \secref{c9:b6}.

Övriga övningar görs med papper och penna. Lösningsförslag finns i bilaga~\ref{app:answers}.

Varje övning är märkt med sin sort. Det finns en \textbf{Kontroll} i det här kapitlet,
övning~\ref{c9:ex:9}.
\end{aside}

\exercise{Rita en flödesplan}{Konstruktion}
\label{c9:ex:1}
Ett program ska läsa ett tal i \code{r16} och tända PB0 om talet är jämnt, annars PB1. Ett tal är
jämnt om bit~0 är 0.
\begin{parts}
  \item Rita flödesplanen.
  \item Vilken instruktion och mask använder du för att ta reda på bit~0? Tips: \secref{c8:b2}.
  \item Vilket villkorligt hopp passar efter den instruktionen?
\end{parts}

\exercise{Vilket hopp?}{Förståelse}
\label{c9:ex:2}
Skriv en jämförelse och ett villkorligt hopp till etiketten \code{yes} för varje villkor. Talen är
utan tecken om inget annat sägs.
\begin{parts}
  \item \code{r16} är lika med \code{r17}.
  \item \code{r16} är inte noll.
  \item \code{r16} är mindre än 100.
  \item \code{r16} är 100 eller större.
  \item \code{r16} är större än 100.
  \item \code{r16} är negativt, när det läses med tecken.
  \item \code{r16} är mindre än \code{r17}, när båda läses med tecken.
\end{parts}

\exercise{Följ programmet}{Räkna för hand}
\label{c9:ex:3}
Vilket värde får \code{PORTB} när \code{r16} börjar med 5, 10, 99, 100 respektive 200?

\begin{asmcode}
    ldi r17, 0
    cpi r16, 10
    brlo small
    cpi r16, 100
    brsh big
    ldi r17, 2
    rjmp done
small:
    ldi r17, 1
    rjmp done
big:
    ldi r17, 3
done:
    out PORTB, r17
\end{asmcode}

\noindent Rita också programmets flödesplan. Vad är programmets uppgift, med en mening?

\exercise{Om, annars}{Program}
\label{c9:ex:4}
\begin{parts}
  \item Skriv ett program som tänder alla lysdioder på port B om \code{r16} är 42, och släcker
    alla annars.
  \item Testa programmet med \code{r16} = 42 och med \code{r16} = 41. Byt värde antingen i
    \code{ldi}-raden eller i Processor Status.
  \item Ta bort \code{rjmp}-instruktionen efter den första grenen. Vad händer nu för
    \code{r16}~=~41, och för \code{r16}~=~42? Förklara.
\end{parts}

\exercise{Hur många varv?}{Räkna för hand}
\label{c9:ex:5}

\begin{asmcode}
    ldi r16, N
loop:
    inc r17
    dec r16
    brne loop
\end{asmcode}

\noindent\code{r17} är 0 från början. Vad är \code{r17} efter loopen när
\begin{parts}
  \item N = 5?
  \item N = 1?
  \item N = 0? Förklara svaret.
  \item Hur många gånger körs \code{brne} när N = 5, och hur många av dem hoppar den?
\end{parts}

\exercise{Summera}{Program}
\label{c9:ex:6}
\begin{parts}
  \item Skriv ett program som räknar ut $1 + 2 + \dots + n$, med $n$ i \code{r16} och summan i
    \code{r17}.
  \item Prova med $n = 15$. Vilket svar ska du få?
  \item Vilket är det största $n$ som ger rätt svar i en byte? Vad får du för $n = 23$, och
    varför?
\end{parts}

\exercise{Räkna cykler}{Räkna för hand}
\label{c9:ex:7}
\begin{parts}
  \item Hur många klockcykler tar fördröjningen i \secref{c9:b4} om startvärdet är 100 i stället
    för 200? Hur lång tid är det vid 16~MHz?
  \item Samma fråga för startvärdet 1.
  \item Samma fråga för startvärdet 0.
\end{parts}

\exercise{Välj konstanter}{Konstruktion}
\label{c9:ex:8}
\begin{parts}
  \item En fördröjning ska vara 40~µs. Hur många varv ska loopen i \secref{c9:b4} gå? Hur nära
    40~µs kommer du, och hur kan du komma exakt?
  \item En fördröjning ska vara 5~ms. Välj OUTER och INNER för loopen i \secref{c9:b5}, först utan
    och sedan med en \code{nop} i den yttre loopen. Hur nära 5~ms kommer du i de två fallen?
  \item Varför går det inte att göra en fördröjning på en sekund med två loopar?
\end{parts}

\exercise{Räkna, mät och förklara}{Kontroll}
\label{c9:ex:9}
\emph{Räkna ut för hand, mät med stoppuret, förklara skillnaden.}

Använd \code{long_delay.asm}, med OUTER = 208 och INNER = 255:

\asmfile{lectures/L09/examples/long_delay.asm}

\begin{parts}
  \item Räkna ut antalet klockcykler från raden \code{delay_start} till raden \code{delay_end},
    med en tabell som den i \secref{c9:b4}: en rad per instruktion, cykler, antal gånger och
    summa. Räkna om till millisekunder.
  \item Mät samma sak i simulatorn med brytpunkter och stoppuret. Stämmer det?
  \item Flytta den första brytpunkten till raden \code{sbi PORTB, 0} och den andra till raden
    \code{cbi PORTB, 0}. Förutsäg först hur mycket mätningen ändras, mät sedan. Förklara
    skillnaden.
  \item Ändra OUTER till 100 och förutsäg den nya tiden med formeln. Mät och jämför.
  \item En kamrat har räknat ut att fördröjningen tar $\text{OUTER} \cdot \text{INNER} \cdot 3 =
    159\,120$ cykler. Vad har kamraten glömt, och hur stort blir felet i procent?
\end{parts}

\exercise{Flerval}{Program}
\label{c9:ex:10}
Ett program ska visa ett mönster på port B beroende på läget i \code{r16}:
\begin{itemize}
  \item läge 0: alla släckta;
  \item läge 1: bara PB0;
  \item läge 2: PB0 och PB1;
  \item alla andra lägen: alla tända, som en felindikering.
\end{itemize}
\begin{parts}
  \item Rita flödesplanen, med en gemensam ruta för \code{PORTB = mönster}.
  \item Skriv programmet och testa alla fyra fallen.
\end{parts}

\exercise{Arbeta om}{Konstruktion}
\label{c9:ex:11}
Programmet nedan fungerar, men det har tre slut och tre kopior av \code{out}. Rita dess flödesplan,
arbeta om den enligt \secref{c9:a8}, och skriv det omarbetade programmet.

\begin{asmcode}
    cpi r16, 1
    breq one
    cpi r16, 2
    breq two
    ldi r17, 0
    out PORTB, r17
end_1:
    rjmp end_1
one:
    ldi r17, 1
    out PORTB, r17
end_2:
    rjmp end_2
two:
    ldi r17, 3
    out PORTB, r17
end_3:
    rjmp end_3
\end{asmcode}

\exercise{Minusgrader}{Förståelse, Fördjupning}
\label{c9:ex:12}
Värmaren i \secref{c9:a5} ska användas utomhus, där temperaturen kan bli negativ. Den lagras som
ett tal med tecken.
\begin{parts}
  \item Vad gör programmet när temperaturen är $-5$ grader? Följ \code{cpi} och \code{brlo} steg
    för steg.
  \item Vilket hopp ska användas i stället, och vilken flagga läser det?
  \item Fungerar det rättade programmet även för temperaturen 25 grader? För 100 grader?
\end{parts}
